\begin{thema}{Δ}
Δίνεται η συνάρτηση $f(x) = (x-1)\ln x - x + 1$, με $x > 0$.
\begin{erwthma}
    \item Να μελετήσετε τη συνάρτηση $f$ ως προς τη μονοτονία και τα ακρότατα.
    \item Να αποδείξετε ότι για κάθε $x>0$ ισχύει $(x-1)\ln x \ge x - 1$. Πότε ισχύει η ισότητα;
    \item Να βρείτε την εξίσωση της εφαπτομένης της $C_f$ η οποία διέρχεται από την αρχή των αξόνων $O(0,0)$.
    \item Θεωρούμε την αρχική συνάρτηση $F$ της $f$ στο $(0, +\infty)$. Να αποδείξετε ότι για κάθε $\alpha, \beta > 1$ με $\alpha \ne \beta$ ισχύει:
    \[ \frac{F(\beta) - F(\alpha)}{\beta - \alpha} > 0 \]
\end{erwthma}
\end{thema}
